\Subhead{6.1.2}{Статистическая проверка гипотез. Критерии Байеса, Неймана--Пирсона и Вальда}

Пусть $C_{ij}$ — цена решения $H_i$, если истинно $H_j$, а $P_j$ — априорная вероятность. Сравнение условных рисков даёт байесовское правило
\begin{equation}
 \Lambda(y)\mathop{\gtrless}_{H_0}^{H_1}
 \frac{(C_{10}-C_{00})P_0}{(C_{01}-C_{11})P_1},
\end{equation}
если разности цен положительны. При равных ценах ошибок и нулевых ценах верных решений выбирают гипотезу с большей апостериорной вероятностью.

Нейман–Пирсон фиксирует вероятность ложной тревоги $P_{\rm FA}=\alpha$ и максимизирует вероятность обнаружения $P_D$. Оптимально сравнивать то же отношение с порогом, выбранным из распределения при $H_0$. Здесь априорные вероятности не нужны. Если статистика дискретна, на границе может понадобиться рандомизация, чтобы точно обеспечить $\alpha$.

В последовательном критерии Вальда независимые наблюдения накапливаются до решения:
\begin{equation}
 L_n=\prod_{j=1}^n\frac{p(y_j|H_1)}{p(y_j|H_0)},\quad
 L_n\geq A\Rightarrow H_1,\quad L_n\leq B\Rightarrow H_0.
\end{equation}
Между границами измерение продолжают. При заданных ошибках $\alpha=P(1|0)$ и $\beta=P(0|1)$ приближённо $A\simeq(1-\beta)/\alpha$, $B\simeq\beta/(1-\alpha)$. Точность границ ограничена перескоками через порог; при коррелированных данных нужны условные плотности, а не произведение маргинальных. Вальд экономит среднее число наблюдений в своей постановке, но время решения случайно. ROC-кривая показывает компромисс $P_D$ и $P_{\rm FA}$ при перестройке порога.
